\documentclass[output=paper]{langscibook}
\ChapterDOI{10.5281/zenodo.22880288}
\title{\textit{Varigomo} – A kind of net fishing at low tide}
\author{Enoch Horai Magum}
\abstract{}
\begin{document}
\maketitle\label{chap:30}
\begin{figure}[H]
    \includegraphics[trim=0 5 0 15,clip,width=.9\linewidth]{figures/30_1.jpg}
    \caption{\textit{Varigomo} -- Net-fishing at low tide}
    \label{fig:30_1}
\end{figure}
\begin{xltabular}{\textwidth}{lXX}
	(\ref{ex:text23-1})&\textit{A varigomo, eove o peho kaku vagana vaa Teapu to paku raara bona kave.}&\textit{Varigomo} is a kind of fishing in the Teop area that we do with a net.\\
	(\ref{ex:text23-2})&\textit{A kave nae he a bebeahu ae a mata rutaa.}&The net used for it is long and has small meshes.\\
	(\ref{ex:text23-3})&\textit{O kaku vagana bona na pakupaku raara teo gohoho ge teo tarahi, ei tea maa hana taataana, sa ei haana namana.}&We do this kind of fishing during the incoming tide or the ebbing tide, this means at places where we can walk, this is not in the deep ocean.\\
	(\ref{ex:text23-4})&\textit{O kaku vagana bona sa vaapenapena haana, beo aba o bero nao, evehee a bero amaa vuri, o kukan koa o sinivi to tahoa rori.}&This kind of fishing does not forbid that many people go, but most of the time it is only three canoes that set off.\\
	(\ref{ex:text23-5})&\textit{O peho sinivi to vaahae nana bona kave mi bona kukan a aba, a buaku vaa tea vavaagunaha kave, a peha he a too horohoro.}&One canoe takes the net with three people on board, two for letting the net down, and one person is the punter.\\
	(\ref{ex:text23-6})&\textit{O meho buo sinivi he mene vahakaa.}&But the two other canoes have to beat the water to chase the fish into the net.\\
	(\ref{ex:text23-7})&\textit{O kaku vagana nae he na paku vonaen:}&The fishing is done as follows:\\
	(\ref{ex:text23-8})&\textit{Beara hua batara, eara toro tara bata naora beo iana an rori.}&While we are paddling along, we look where the fish eat.\\
	(\ref{ex:text23-9})&\textit{Bea too horohoro sue voen na mataa nana, eara repaa hua vahuus o pata iana bari.}&When the punter says, it is good, we paddle close to that school of fish.\\
	(\ref{ex:text23-10})&\textit{A peha aba repaa gunaha me bono peho punete kave.}&Then one person goes down with one edge of the net.\\
	(\ref{ex:text23-11})&\textit{O sinivi repaa hua taviri bono pata iana, bea meha aba rurun gunaha bata ni nana bona kave.}&The canoe then paddles around the school of fish, while the other man is dropping the net along the way.\\
	(\ref{ex:text23-12})&\textit{O meho buo sinivi he paa taneo vaevuru tea vahakaa, o meho tea papa matau, o meho he tea papa keruka.}&The other two canoes have already started to beat the water, one on the right side, the other one on the left side.\\
	(\ref{ex:text23-13})&\textit{Pasi vatete maari bono pata iana, ore vatete hovo komana kave.}&They will chase the school of fish, and they chase them into the net.\\
	(\ref{ex:text23-14})&\textit{Bea bua too vavaagunaha kave rahi vatatana bata ni rori bono buo punete kave, eori repaa asun bona iana en komana kave.}&When the two net-droppers pull the two edges of the net together, they kill the fish inside the net.\\
	(\ref{ex:text23-15})&\textit{Amaa aba vai to sa paa hana hau tea maa mata kave, kahi tahitahi ri vuara bona makaree.}&Those that have not been caught in the meshes of the net, we will spear with the fishing spear.\\
	(\ref{ex:text23-16})&\textit{Amaa aba vai he to paa hanau teo mata kave na asun ri raara bono popoura.}&But those that hang in the meshes of the net, we kill them with the baton.\\
	(\ref{ex:text23-17})&\textit{Be kavara, eara repaa vaahae vahabana a kave, eara repaa  nao pete teo meho hum vai to anan bata rio iana.}&After that we take the net again on board and go to another place where the fish eat.\\
	(\ref{ex:text23-18})&\textit{A vaakavara vagana na tei bata koana tea masi ge beo gohoho beera oha maa, eara repaa voosu vaha vaan.}&The end of the fishing trip is when it is low tide and when the incoming tide is too big; then we go home, back to the village.\\
\end{xltabular}
\begin{figure}[H]
    \includegraphics[trim=0 0 0 0,clip,width=\linewidth]{figures/30_2.png}
    \caption{\textit{Makaree} -- Fishing spear}
    \label{fig:30_2}
\end{figure}
\begin{figure}[H]
    \includegraphics[trim=0 0 0 0,clip,width=.9\linewidth]{figures/30_3.png}
    \caption{\textit{Popoura} -- Baton}
    \label{fig:30_3}
\end{figure}
\end{document}